% Про необходимое условие теоремы
\begin{example}
	Требование непрерывности частных производных существенно. Это показывает следующая функция:
	\[
		f(x, y) = \sqrt{|xy|}
	\]
	Покажем, что у этой непрерывной в нуле функции есть частные производные в этой же точке, но при этом они разрывны и сама функция не дифференцируема. Начнём с частных производных:
	\[
		\pd{f}{x}(0, 0) = \liml_{\Delta x \to 0} \frac{f(\Delta x, 0) - f(0, 0)}{\Delta x} = 0 = \pd{f}{y}(0, 0)
	\]
	Сами функции частных производных имеют вид
	\[
		\pd{f}{x}(x, y) = \sqrt{|y|} \cdot \frac{1}{2\sqrt{|x|}} \cdot \sgn{x} \xrightarrow[x \to 0]{} \infty
	\]
	Значит, у нас присутствует разрыв частных производных в нуле.
\end{example}

\begin{example}
	Теорема в обратную сторону, вообще говоря, неверна. Рассмотрим функцию $f(x, y)$ следующего вида:
	\[
		f(x, y) = \sqrt[3]{x^2 y^2}
	\]
	Частные производные в точке $(0, 0)$, очевидно, равны нулю и при этом выполнено следующее:
	\[
		\Delta f(0, 0) = \sqrt[3]{(\Delta x)^2 (\Delta y)^2} = r^{4/3} \sqrt[3]{\cos^2 \phi \sin^2 \phi} \le r^{4/3}
	\]
	То есть $\Delta f = o(r),\ r \to 0$. Таким образом функция дифференцируема, хоть ее частные производные в нуле не непрерывны. Таким образом мы показали отсутствие необходимости условий теоремы для дифференцируемости.
\end{example}

\begin{theorem} (Дифференцируемость сложной функции)
	Пусть $f$ дифференцируема в точке $\vv{y}_0 \in \R^m$, $\forall j \in \range{m}\ g_j$ дифференцируема в точке $\vv{x}_0 \in \R^n$, причём $g_j(\vv{x}_0) = y_{j, 0}$, если $\vv{y}_0 = (y_{1, 0}, \ldots, y_{m, 0})$. Тогда сложная функция $h(\vv{x}) = f(g_1(\vv{x}), \ldots, g_m(\vv{x}))$ дифференцируема в точке $\vv{x}_0$, причём
	\[
		\grad h(\vv{x}_0) = \grad f(\vv{y}_0) \begin{pmatrix}
			\grad g_1(\vv{x_0})\\
			\vdots\\
			\grad g_n(\vv{x_0})
		\end{pmatrix}
	\]
	где $\grad g_i(\vv{x_0})$ в матрице - её строки, а не элементы.
\end{theorem}

\begin{proof}
	Чтобы доказать дифференцируемость функции $h$, нужно расписать её приращение:
	\[
		\Delta h(\vv{x}_0) = f(g_1(\vv{x}_0 + \vv{\Delta x}), \ldots, g_m(\vv{x}_0 + \vv{\Delta x})) - f(g_1(\vv{x}_0), \ldots, g_m(\vv{x}_0))
	\]
	Положим по определению $\Delta y_j := g_j(\vv{x}_0 + \vv{\Delta x}) - g_j(\vv{x}_0)$. Тогда
	\[
		\forall j \in \range{m}\ g_j(\vv{x}_0 + \vv{\Delta x}) = y_{j, 0} + \Delta y_j
	\]
	То есть приращение можно переписать в другом виде
	\[
		\Delta h(\vv{x}_0) = f(\vv{y}_0 + \vv{\Delta y}) - f(\vv{y}_0) = \frac{\vdelta f}{\vdelta y_1} (\vv{y}_0) \Delta y_1 + \ldots + \frac{\vdelta f}{\vdelta y_m} (\vv{y}_0) \Delta y_m + o(|\vv{\Delta y}|), |\vv{\Delta y}| \to \vv{0}
	\]
	Продолжим цепочку преобразований приращения:
	\begin{multline*}
		\Delta h(\vv{x}_0) = \suml_{j = 1}^m \pd{f}{y_j}(\vv{y}_0) (g_j(\vv{x}_0 + \vv{\Delta x}) - g_j(\vv{x}_0)) + o(|\vv{\Delta y}|) =
		\\
		\suml_{j = 1}^m \frac{\vdelta f}{\vdelta y_j} (\vv{y}_0) \left(\suml_{k = 1}^n \frac{\vdelta g_j}{\vdelta x_k} (\vv{x}_0) \Delta x_k + o_j(|\vv{\Delta x}|)\right) + o(|\vv{\Delta y}|) =
		\\
		\suml_{k = 1}^n \left(\suml_{j = 1}^m \frac{\vdelta f}{\vdelta y_i} (\vv{y}_0) \cdot \frac{\vdelta g_j}{\vdelta x_k} (\vv{x}_0)\right) \Delta x_k + \suml_{j = 1}^m \frac{\vdelta f}{\vdelta y_i} (\vv{y}_0) \cdot o_j(|\vv{\Delta x}|) + o(|\vv{\Delta y}|)
	\end{multline*}
	Если мы докажем, что всё кроме первой суммы  - это $o(|\vv{\Delta x}|)$, то теорема будет доказана. Вторая сумма - конечная сумма выражений вида $o(|\vv{\Delta x}|)$, она вся $o(|\vv{\Delta x}|)$. Осталось показать, что $o(|\vv{\Delta y}|) = o(|\vv{\Delta x}|)$. То есть, что $\frac{o(|\vv{\Delta y}|)}{|\vv{\Delta x}|} \xrightarrow[\vv{\Delta x} \to \vv{0}]{} 0$. $ o(|\vv{\Delta y}|) = \alpha(\vv{\Delta y})|\vv{\Delta y}|$,  где
	$\alpha( \vv{\Delta y}) \to \vv{0}$, при
	 $\ \vv{\Delta y} \to \vv{0}$, а т.к. $\vv{\Delta y} \to \vv{0}$, при $ \vv{\Delta x} \to \vv{0}$, то $ o(|\vv{\Delta y}|) = \beta(\vv{\Delta x})|\vv{\Delta y}|$,  где
	 $\beta( \vv{\Delta x}) \to \vv{0}$, при
	 $\ \vv{\Delta x} \to \vv{0}$. \\
	 Значит (здесь вообще второй переход должен быть к сумме по всем $g_j$, но мы для краткости покажем что для одной есть стремление к нулю, а значит и для всех)
	 \[
	 	o(|\vv{\Delta y}|) = \beta(\vv{\Delta x})|\vv{\Delta y}| \leq \beta(\vv{\Delta x}) \cdot | \trbr{\grad g(\vv{x_0}), \vv{\Delta x}} + o(\vv{\Delta x}))| \le \beta(\vv{\Delta x}) \cdot | \grad g(\vv{x_0})||\vv{\Delta x}| + o(\vv{\Delta x}) = o(\vv{\Delta x})
	 \]
	 Точнее 
	 Последнее равенство следует из того, что $\beta \cdot \grad$ всё еще стремится к нулю, раз $\beta$ стремится к нулю, а $\grad$ - фиксированный, значит вся сумма - $o$-малое просто по определению. 
\end{proof}


\begin{corollary}
	\textcolor{red}{Отсутствует в нынешнем изложении, но определённо удобно для понимания.}
	Сумма, разность и произведение функций многих переменных - дифференцируема в точке $\vv{x}_0$, если изначальные функции тоже были дифференцируемы в точке $\vv{x}_0$
\end{corollary}

\begin{corollary} (Инвариантность формы первого дифференциала)
	Пусть $f$ дифференцируема в точке $\vv{y}_0$. Тогда формула
	\[
		df(\vv{y}_0) = \suml_{j = 1}^m \frac{\vdelta f}{\vdelta y_j} (\vv{y}_0) dy_j
	\]
	справедлива как в случае, когда $y_1, \ldots, y_m$ - независимые переменные, то и в случае, когда они являются дифференцируемыми функциями от $x_1, \ldots, x_n$
\end{corollary}

\begin{proof}
	Пусть $f$ зависит от $y_j = g_j(x_1, \ldots, x_n)$. Тогда $dh(\vv{x}_0)$, где $h = f(g_1, \ldots, g_m)$, имеет вид
	\[
		dh(\vv{x}_0) = \suml_{k = 1}^n \frac{\vdelta h}{\vdelta x_k} (\vv{x}_0) dx_k = \suml_{k = 1}^n \left(\suml_{j = 1}^m \frac{\vdelta f}{\vdelta y_j} (\vv{y}_0) \frac{\vdelta g_j}{\vdelta x_k} (\vv{x}_0)\right) dx_k =
		\\
		\suml_{j = 1}^m \frac{\vdelta f}{\vdelta y_j} (\vv{y}_0) \underbrace{\left(\suml_{k = 1}^n \frac{\vdelta g_j}{\vdelta x_k} (\vv{x}_0) dx_k\right)}_{dy_j}
	\]
\end{proof}

\subsection{Геометрический смысл градиента и дифференцируемости}

\begin{definition}
	\textit{Параметрически заданной поверхностью} называется множество точек $(x(u,v), y(u,v), z(u,v))\in\R^3$, где $(u,v ) \in\overline{D} \subset \R^2$, $x, y, z$ - непрерывные, $\overline{D}$ - некоторая замкнутая область (замыкание области).
\end{definition}

\begin{definition}
	Плоскость, проходящая через точку $M_0$, называется \textit{касательной плоскостью} к заданной поверхности в данной точке $M_0$, если угол между секущей проходящей через $M_0$ и другую точку $M_1$ поверхности, и плоскостью стремится к нулю при $M_1 \to M_0$.
\end{definition}

\begin{note}
	График непрерывной функции $z = f(x,y) \in \overline{D} \subset \R^2 $, является поверхностью.
\end{note}

% Здесь нужна картинка с осями и плоскостями. Пнуть меня для фото

\begin{theorem}
	Если $f:\overline{D}\to\R, \overline{D}$ - замкнутая область в $\R^2$, дифференцируема в точке $(x_0, y_0) \in \overline{D}$, то касательная плоскость к графику функции $f$ в точке $M_0(x_0, y_0, f(x_0, y_0))$ существует и задаётся уравнением
	\[
		z - z_0 = \pd{f}{x} (x_0, y_0) (x - x_0) + \pd{f}{y} (x_0, y_0) (y - y_0)
	\]
	где $z_0:=f(x_0,y_0)$.
\end{theorem}

\begin{proof}
	Выберем точку $M_1(x, y, f(x, y))$. Из дифференцируемости следует, что $f$ определена в некоторой окрестности точки $(x_0, y_0)$. Стало быть, $M_1$ в окрестности тоже определена. Запишем уравнение вектора $\vv{M_0 M_1}$:
	\[
		\vv{M_0M_1} \leftrightarrow (x - x_0, y - y_0, f(x, y) - f(x_0, y_0))
	\]
	Нормальный вектор $\vv{n}$ к плоскости, заданной уравнением в теореме, имеет вид
	\[
		\vv{n} \leftrightarrow \left(\pd{f}{x} (x_0, y_0), \pd{f}{y} (x_0, y_0), -1\right)
	\]
	Нам нужно показать, что угол между этими вектора стремится к $\frac{\pi}{2}$, а для этого достаточно показать, что его косинус стремится к нулю, чем мы и займёмся:
	\begin{multline*}
		|\cos \angle (\vv{M_0M_1}, \vv{n})| = \frac{|\trbr{\vv{M_0M_1},\vv{n}}|}{|\vv{M_0M_1}|\cdot |\vv{n}|} = \\ = \frac{|\pd{f}{x} (x_0, y_0)(x - x_0) + \pd{f}{y} (x_0, y_0)(y - y_0) - f(x, y) + f(x_0, y_0)|}{\sqrt{(\pd{f}{x} (x_0, y_0)^2+ \pd{f}{y} (x_0, y_0)^2+1}\cdot \sqrt{(x-x_0)^2 + (y-y_0)^2 + (f(x,y)-z_0)^2}}
	\end{multline*}
	Также распишем дифференцируемость функции в $(x_0, y_0)$:
	\[
		z - f(x_0, y_0) = \frac{\vdelta f}{\vdelta x}(x_0, y_0)(x - x_0) + \frac{\vdelta f}{\vdelta y}(x_0, y_0)(y - y_0) + o\big(|(x - x_0, y - y_0)|\big)
	\]
	Итого, подставив, получаем:
	\begin{multline*}
		|\cos \angle (\vv{M_0M_1}, \vv{n})| = \frac{o(|(x - x_0, y - y_0)|)}{\sqrt{(\pd{f}{x} (x_0, y_0)^2+ \pd{f}{y} (x_0, y_0)^2+1}\cdot \sqrt{(x-x_0)^2 + (y-y_0)^2 + (f(x,y)-z_0)^2}} \leq \\ \leq  \frac{o(|(x - x_0, y - y_0)|)}{\sqrt{1} \cdot \sqrt{(x - x_0)^2 + (y - y_0)^2}} \xrightarrow[(x, y) \to (x_0, y_0)]{} 0
	\end{multline*}
\end{proof}

\begin{definition}
	\textit{Производной функции $f$ в точке $\vv{x}_0 \in \R^n$ по направлению} $\vv{l}$ называется предел (если он существует):
	\[
		\pd{f}{\vv{l}} (\vv{x}_0) := \liml_{t \to 0+} \frac{f(\vv{x}_0 + t\vv{l}) - f(\vv{x}_0)}{t}
	\]
\end{definition}

\begin{note}
	В разных книгах производную по направлению определяют по-разному. Например, могут убрать стремление $t$ к 0 лишь с положительной стороны, могут добавить модуль $\vv{l}$ в знаменатель или потребовать, что $\vv{l}$ имеет единичную длину. Поэтому, если нужно прочитать доказательство, использующее производную по направлению, то стоит уточнить, что именно под этим понятием подразумевает автор.
\end{note}

\begin{theorem}
	Если $f$ дифференцируема в $\vv{x}_0\in \R^n$, то она имеет производную по любому направлению $\vv{l}$, причём
	\[
		\pd{f}{\vv{l}} (\vv{x}_0) = \trbr{\grad f(\vv{x}_0), \vv{l}}
	\]
\end{theorem}

\begin{proof}
	Если $\vv{l} = \vv{0}$, то и скалярное произведение выродится в ноль, и производная по направлению(чисто по определению). Теперь случай $\vv{l} \neq \vv{0}$. Распишем дробь под пределом:
	\begin{multline*}
		\pd{f}{\vv{l}}(\vv{x}_0) = \liml_{t \to 0^+}\frac{f(\vv{x}_0 + t\vv{l}) - f(\vv{x}_0)}{t} = \liml_{t \to 0^+}\frac{\trbr{\grad f(\vv{x}_0), t\vv{l}} + o(|t\vv{l}|)}{t} = \\
		= \liml_{t \to 0^+} \left(\trbr{\grad f(\vv{x}_0), \vv{l}} + |\vv{l}| \cdot \frac{o(t)}{t}\right) = \trbr{\grad f(\vv{x}_0), \vv{l}}
	\end{multline*}
\end{proof}

\begin{theorem}
	(Геометрический смысл градиента)
	Если $f$ дифференцируема в точке $\vv{x}_0\in\R^n$ и $\grad f(\vv{x}_0) \neq \vv{0}$, то производная по направлению $\vv{l},\ |\vv{l}| = 1$
	\begin{itemize}
		\item Максимальна при $\vv{l} = \frac{\grad f(\vv{x}_0)}{|\grad f(\vv{x}_0)|}$
		
		\item Минимальна при $\vv{l} = -\frac{\grad f(\vv{x}_0)}{|\grad f(\vv{x}_0)|}$
	\end{itemize}
\end{theorem}

\begin{anote}
	То есть по сути данное следствие указывает, что изменение функции максимально по направлению градиента в данной точке и минимально в обратную сторону. Читай, график наиболее крутой$\backslash$пологий в эти стороны.
\end{anote}

\begin{proof}
	Распишем модуль производной по направлению как скалярное произведение:
	\[
		\left|\pd{f}{\vv{l}} (\vv{x}_0)\right| = \left|\trbr{\grad f(\vv{x}_0), \vv{l}}\right| \le |\grad f(\vv{x}_0)|\cdot|\vv{l}|=|\grad f(\vv{x}_0)|
	\]
	Равенство достигается, если и только если $\grad f(\vv{x}_0)$ и $\vv{l}$ коллинеарны(смотри неравенство Коши-Буняковского-Шварца). Тогда легко видеть, что следующие значения и только они соответствуют искомым максимуму и минимуму: 
	\[
		\vv{l} = \pm \frac{\grad f(\vv{x}_0)}{|\grad f(\vv{x}_0)|}
	\]
\end{proof}